\documentclass[11pt]{article}
\usepackage[margin=2.5cm]{geometry}
\usepackage{booktabs}
\usepackage{caption}

\title{Emergence of the number line across pretraining checkpoints}
\author{}
\date{}

\begin{document}
\maketitle

\noindent For each checkpoint, every four-digit number of a fixed set of 2,000 is represented by the hidden state at
the final \texttt{=} of its copy prompt ($\sim$84\% depth). For each number we take its 10 nearest numbers
(Euclidean distance) and report the share that are within $\pm 50$ in value, share its leading digit, or share
its last digit, together with the model's first-token accuracy on the copy task (300 prompts).

\begin{table}[ht]
\centering
\small
\begin{tabular}{rcccc}
\toprule
 & \multicolumn{3}{c}{Share of the 10 nearest neighbours} & Copy task \\
\cmidrule(lr){2-4}
Step & within $\pm 50$ & same leading digit & same last digit & first-token acc. \\
\midrule
0 & 1\% & 12\% & 11\% & 0\% \\
64 & 2\% & 12\% & 11\% & 0\% \\
256 & 2\% & 13\% & 11\% & 1\% \\
512 & 2\% & 15\% & 13\% & 2\% \\
1{,}000 & 3\% & 16\% & 12\% & 1\% \\
2{,}000 & 14\% & 36\% & 16\% & 12\% \\
4{,}000 & 55\% & 95\% & 13\% & 78\% \\
8{,}000 & 57\% & 97\% & 14\% & 93\% \\
13{,}000 & 62\% & 99\% & 17\% & 100\% \\
23{,}000 & 69\% & 99\% & 20\% & 100\% \\
33{,}000 & 68\% & 99\% & 20\% & 100\% \\
53{,}000 & 68\% & 98\% & 21\% & 100\% \\
73{,}000 & 65\% & 94\% & 26\% & 99\% \\
103{,}000 & 67\% & 97\% & 20\% & 100\% \\
123{,}000 & 65\% & 97\% & 19\% & 100\% \\
143{,}000 & 65\% & 97\% & 19\% & 100\% \\
\midrule
Chance & 1\% & 11\% & 11\% & -- \\
\bottomrule
\end{tabular}
\caption{Pythia-1.4B: neighbourhood composition of four-digit numbers ($\sim$84\% depth) and copy accuracy across training. The number line appears between steps 2{,}000 and 4{,}000, together with the copy skill.}
\label{tab:numline_pythia}
\end{table}

\begin{table}[ht]
\centering
\small
\begin{tabular}{rcccc}
\toprule
 & \multicolumn{3}{c}{Share of the 10 nearest neighbours} & Copy task \\
\cmidrule(lr){2-4}
Step & within $\pm 50$ & same leading digit & same last digit & first-token acc. \\
\midrule
0 & 1\% & 12\% & 15\% & 0\% \\
300 & 1\% & 11\% & 20\% & 0\% \\
10{,}000 & 49\% & 67\% & 11\% & 84\% \\
20{,}000 & 73\% & 88\% & 10\% & 99\% \\
40{,}000 & 86\% & 95\% & 10\% & 100\% \\
70{,}000 & 76\% & 89\% & 11\% & 100\% \\
140{,}000 & 70\% & 84\% & 12\% & 100\% \\
250{,}000 & 59\% & 79\% & 29\% & 100\% \\
420{,}000 & 64\% & 82\% & 20\% & 100\% \\
710{,}000 & 67\% & 83\% & 19\% & 100\% \\
1{,}170{,}000 & 61\% & 81\% & 33\% & 98\% \\
1{,}907{,}359 & 54\% & 77\% & 39\% & 99\% \\
\midrule
Chance & 1\% & 11\% & 11\% & -- \\
\bottomrule
\end{tabular}
\caption{OLMo-2-1B: as Table~\ref{tab:numline_pythia}. The number line appears by step 10{,}000, peaks at step 40{,}000 and weakens later, while the last-digit share rises.}
\label{tab:numline_olmo}
\end{table}

\end{document}
